\section{Experiments}
\label{sec:experiments}

\subsection{Setup}
\label{sec:exp:setup}

Dataset scale evolved through the project: a 1{,}024-clip $\times$ 128-shape pilot $\to$ Stage 1's
full 20{,}946-clip library at one neutral shape $\to$ a 20{,}946-clip $\times$ 2-shape
intermediate set (population-median male/female bodies, used to stress-test MoE without waiting on
full 128-shape data) $\to$ a 150-clip $\times$ 128-shape reduced-scale testbed adopted for cost
reasons as the standard architecture/reward ablation corpus
(Sections~\ref{sec:exp:architecture}--\ref{sec:exp:reward-source}) $\to$ full
20{,}946-clip $\times$ 128-shape training via \texttt{GlobalClipPool}
(Section~\ref{sec:method:infra}) for the final run.

\textbf{Metrics:} $\mathrm{dp\_error}$ (DOF-angle, shape-invariant — primary metric, full
definition and justification in Section~\ref{sec:evaluation}), $\mathrm{gt\_error}$/
$\mathrm{gr\_error}$ (world-space position/rotation, scale-confounded across shapes, kept for
legacy/reward-design tracking), success rate, held-out clip evaluation.

\subsection{Main Results}
\label{sec:exp:main-results}

Stage 1 neutral pretrain reaches 84.9\% success at epoch 20{,}200 ($\sim$174h, 6 GPUs), with an
8.7\% persistent-failure tail (Section~\ref{sec:method:curriculum}). For Stage 2 shape-transfer,
a residual-adapter lineage (frozen backbone + adapter, v1--v6) plateaued at 78--82\% success
across every variant tried and was abandoned in favor of a from-scratch Stage 2 trunk, still in
progress \todo{final Stage 2 success numbers pending}. Stage 2 results should be read alongside the
shape-consistency metrics of Section~\ref{sec:evaluation}, not success rate in isolation.

\subsection{Architecture Evolution: Flat MLP $\to$ MoE $\to$ Self-Attention}
\label{sec:exp:architecture}

Presented as a narrative of iterative design — the actual chronological order tried, each step
motivated by the failure mode of the last — validated at reduced scale (150-clip / 1{,}024-clip
corpora) rather than full scale, for cost reasons.

\begin{enumerate}
  \item \textbf{Flat-MLP baseline.} Wide MLP over flat-concatenated observations (history/lookahead
  frames concatenated, not tokenized). Establishes a $\sim$75--80\% success-rate ceiling on the
  hardest-clip subset.
  \item \textbf{Mixture-of-Experts (MoE).} Hypothesis: distinct motion classes
  (crawl/kneel/squat/turning) benefit from routed specialist sub-policies. \textbf{Result: null} —
  a capacity-matched (parameter-matched, $\sim$95\%) wide-MLP control caught up to the MoE-stable
  variant on all metrics; whether routing helps remains an open question.
  \item \textbf{Self-attention temporal encoder.} Hypothesis: tokenizing per-frame history/
  lookahead and attending (reusing the existing \texttt{Transformer} module) helps structure
  learning over the dilated frame sequence. \textbf{Result: null} — no separation from flat-concat
  at matched training steps.
\end{enumerate}

Both alternatives were null/inconclusive at the scale tested; we retain the simpler flat-MLP
architecture, carried forward as a bundled recipe into the one full-scale run rather than
re-validated at full scale.

\subsection{Morphology Conditioning Ablation: Raw Betas vs.\ Physics-Derived Features}
\label{sec:exp:morphology}

Pilot-scale comparison ($1{,}024$-clip $\times$ 128-shape corpus): raw SMPL betas
(21{,}400 epochs) vs.\ a 15-dim $z$-scored physics-derived feature vector (mass, height,
limb-length-style scalars; 17{,}200 epochs). Persistent-failure-clip counts (177 vs.\ 167) are
\textbf{not statistically significant}; qualitative visual review favored physics features (5/8
vs.\ 0/8 clips judged visually improved). Raw betas are kept as the simpler default for the
full-scale run — the physics-feature signal was not strong enough to justify the added
feature-engineering surface at $20{,}946 \times 128$ scale.

\subsection{AMASS/HUMOS Reward Source Design}
\label{sec:exp:reward-source}
\todo{In progress, not yet resolved.}

Motivation: AMASS-canonical DOF-space targets are exactly shape-invariant
(Section~\ref{sec:method:data}), so training on them alone gives zero shape-conditioning gradient;
HUMOS is the only channel whose target genuinely differs by shape. Iterative design, same framing
as Sections~\ref{sec:exp:architecture}--\ref{sec:exp:morphology}:
\begin{enumerate}
  \item DOF-space-only bundle on the canonical corpus — clean $\mathrm{dp\_error}$ convergence,
  world-space metrics flat.
  \item Episode-level source-switching (hard-mask reward terms by AMASS-vs-HUMOS per episode) —
  world-space metrics improve over (1) but underperform the unmasked world-space baseline at
  matched steps (63\% vs.\ 83\% success), likely because masking halves how often the world-space
  reward fires per rollout.
  \item Combined, unmasked, full-factor reward on the mixed corpus with per-source \emph{weights}
  (not term masks) — outcome pending.
\end{enumerate}

\subsection{Reward/Termination Ablation}
\label{sec:exp:reward-term-ablation}

Discover-relaxed vs.\ a stricter baseline configuration (full termination + effort/smoothness/
contact-timing terms), validated at reduced scale for the same cost reason as
Sections~\ref{sec:exp:architecture}--\ref{sec:exp:morphology}, then carried forward as-is to full
scale by explicit decision.

\subsection{Failure Analysis}
\label{sec:exp:failure}

Persistent-failure motion clusters show substantial overlap across single-shape and multi-shape
settings: 97.1\% overlap between the full-scale-scratch hardest-clip scoreboard and the
single-shape neutral persistent-failure set — the plateau is an intrinsic motion-difficulty
ceiling, not shape- or reward-specific. Shape-failure correlation (Pearson $r \approx -0.09$ to
$0.003$ between shape extremity and failure) is reported here in full; the headline framing of the
same result appears in Section~\ref{sec:evaluation}.

\subsection{Cross-Shape Generalization}
\label{sec:exp:cross-shape}

Held-out beta interpolation/extrapolation (experiment E7) is the strongest form of
Section~\ref{sec:evaluation}'s shape-consistency claim, extending it past the 128 training shapes
to unseen interpolated/extrapolated shapes. \todo{Blocked on an \texttt{infer.py} held-out-beta
pipeline bug; deferred in favor of reusing already-existing held-out-shape data once a Stage 2
checkpoint is ready.}
